\subsubsection{\texorpdfstring{$\mathbf{PW}$}{PW} Type safety: unitary slice addressability and null value (NULL)}\label{subsec:xi-pw-type-safety-null}
The \(\Xi_{\Delta,L}\) ​​zero-point audit scope in this chapter is based on the fixed geometry of the unitary slice (\(|u|=1\)) (note \ref{rem:xi-critical-line-unit-circle}), and uses the spectral crossing on the \(\theta\)-axis as a reviewable interface (proposition \ref{prop:xi-unitary-slice-spectral-crossing} and note \ref{rem:xi-unitary-slice-operator-audit}).
On the other hand, in the semantics of projected ontology, the concept starts with \emph{address} (Section \ref{sec:recursive-addressing}): without an address, it is read as undefined (null value), and with an address, it is read as a verifiable projection/certificate.
This section compiles "the point deviating from \(\Re(s)=\tfrac12\) ​​is read as NULL (no address)" into a \(\mathbf{PW}\) type safety (soundness) proposition within the system, and explicitly distinguishes it from the completeness bridge required by RH.

\paragraph{"Visible domain" and "projection gate" in unit circle area phase arithmetic}
Appendix \ref{app:unit-circle-phase-gate} ​​The parameter-free compactification door given
\(\upsilon(x)=\pi^{-1}\arctan(x)\in(-\tfrac12,\tfrac12)\)
and its unit circle equivalent embedding
\(\iota(x)=\frac{1+\mathrm{i}x}{1-\mathrm{i}x}\in\TT\)
On the projective compactification \(\widehat{\RR}:=\RR\cup\{\infty\}\) ​​it can be extended to the homeomorphism \(\iota:\widehat{\RR}\to\TT\) (\S\ref{subsec:unit-circle-phase-lift}).
Therefore, \emph{only at the coordinate gate level} does not necessarily have null values: the compactification gate itself is a fully defined reversible mapping, and any irreversibility comes from the subsequent "partial observation/projection readout" (Appendix \S\ref{subsec:unit-circle-phase-information}) and the \(\mathsf{NULL}\) ​​convention of address semantics (see Section \ref{sec:spg}).
\par
The key to this section is: when the verification interface of \(\Xi\) ​​is fixed as a unitary slice, the verifiable continuous parameter axis is further compressed into a \emph{visible slice};
Even if the object that deviates from this slice can be written formally at the parsing level, it has no address under the closure of this paper, and its readout is degenerated to \(\mathsf{NULL}\) ​​according to the null value semantics.
Under the scope of the standard group embryo, the visible slice can be regarded as a standard section, and the degrees of freedom that deviate from the slice belong to the standard fiber; if there is no external recording axis, it will be processed as \(\mathsf{NULL}\) ​​in the type layer (see section \ref{subsec:pom-gauge-groupoid}).

\paragraph{Phase type parameter domain (unitary slice)}
remember
\[
\mathsf{Phase}:=\{u\in\CC:\ |u|=1\}=\{e^{i\theta}:\ \theta\in\RR\}.
\]
For any \(L>1\), we get by defining \(u_L(s)=L^{2s-1}\)
\[
\Re(s)=\tfrac12\quad\Longleftrightarrow\quad |u_L(s)|=1\quad\Longleftrightarrow\quad u_L(s)\in \mathsf{Phase}
\]
(Note \ref{rem:xi-critical-line-unit-circle}).

\begin{definition}[Readout under visible domain constraints and \texorpdfstring{$\mathsf{NULL}$}{NULL}]\label{def:xi-visible-read-null}
Given a state space \(\Omega\), an observation map \(\operatorname{obs}:\Omega\rightharpoonup X\) (partial mapping is allowed), and a visible domain \(X_{\mathrm{vis}}\subseteq X\).
And fix a closure/checker \(\mathrm{Chk}\) ​​(which in the context of this paper can be taken as the combination of "\(\mathbf{PW}\) regularized closure + unitary-slice audit interface": its output is \(\mathsf{pass}/\mathsf{fail}\)).
Define visible readouts with visible domain constraints
\[
\mathrm{read}(\omega)=
\begin{cases}
\operatorname{obs}(\omega), &
\omega\in\mathrm{Dom}(\operatorname{obs}),\ \operatorname{obs}(\omega)\in X_{\mathrm{vis}},\ \mathrm{Chk}(\omega)=\mathsf{pass},\\
\mathsf{NULL}, & \text{otherwise}.
\end{cases}
\]
Among them, \(\mathrm{read}(\omega)=\mathsf{NULL}\) ​​means that the readout is not defined (not addressable).
This "otherwise" can be decomposed into three types of structural sources (all occurring at the \emph{referential/type level} rather than the numerical level):
\begin{enumerate}
  \item \textbf{Semantic null value (missing address):} The address precedes the value; if the address is not given, the corresponding conceptual value is undefined, recorded as \(\mathsf{NULL}\), and is not equivalent to the numerical value \(0\) (see section \ref{sec:spg}; see also section \ref{sec:recursive-addressing}).
  \item \textbf{Protocol null (forbidden by the type system):} The object's interpretation requires degrees of freedom beyond \(X_{\mathrm{vis}}\), or requires axes explicitly prohibited by the interface protocol (typically: radial mode of \(|u|\neq 1\) is required, while the visible domain only allows unitary slices), and thus is not typeable under the given closure/checker and is rejected.
  When \item \textbf{collision null value (insufficient side information):} When \(\operatorname{obs}\) is many-to-one compression and wants to improve the readout to (approximately) reversible with a given confidence scope, an external recording axis must be used to make the extended readout an injective; if the recording axis budget is insufficient, different pre-images cannot be uniquely reconstructed within the allowed error/confidence sense of the protocol, so the readout is still equivalently degenerated to \(\mathsf{NULL}\) (see proposition \ref{prop:xi-null-three-way}).
\end{enumerate}
In the \(\Xi\)-interface in this section, \(\Omega=\CC\), \(\operatorname{obs}(s)=u_L(s)=L^{2s-1}\) are available, and let \(X_{\mathrm{vis}}=\mathsf{Phase}\);
Therefore, \(\Re(s)\neq\tfrac12\Rightarrow \operatorname{obs}(s)\notin X_{\mathrm{vis}}\), its readout directly degenerates into \(\mathsf{NULL}\) by definition.
\end{definition}

\begin{definition}[Visible symmetry (read automorphism)]\label{def:xi-read-visible-symmetry}
The notation that defines \ref{def:xi-visible-read-null} ​​is used.
A bijection \(g:\Omega\to\Omega\) ​​is called a \textbf{visible symmetry} if it exists for any \(\omega\in\Omega\)
\[
\mathrm{read}(g\omega)=\mathrm{read}(\omega).
\]
Let the set of all visible symmetries be \(\mathrm{Aut}_{\mathrm{read}}(\Omega)\).
\end{definition}

\begin{proposition}[Visible symmetrical orbitals and readout fibers]\label{prop:xi-read-fiber-orbits}
The notation that defines \ref{def:xi-read-visible-symmetry} ​​is used.
Then the action orbit of \(\mathrm{Aut}_{\mathrm{read}}(\Omega)\) ​​on \(\Omega\) is exactly the fiber of \(\mathrm{read}\): for any \(\omega\in\Omega\)
\[
\mathrm{Aut}_{\mathrm{read}}(\Omega)\cdot\omega
=\{\omega'\in\Omega:\ \mathrm{read}(\omega')=\mathrm{read}(\omega)\}.
\]
In particular, under the projective ontology semantics, "symmetry" can be understood as reading out the automorphic degrees of freedom within the fiber; this is consistent with the canonical group embryonic scope of \S\ref{subsec:pom-gauge-groupoid}.
\end{proposition}

\begin{proof}
By definition \ref{def:xi-read-visible-symmetry}, for any \(g\in\mathrm{Aut}_{\mathrm{read}}(\Omega)\) ​​and any \(\omega\in\Omega\), there is \(\mathrm{read}(g\omega)=\mathrm{read}(\omega)\), so the orbit \(\mathrm{Aut}_{\mathrm{read}}(\Omega)\cdot\omega\subseteq \mathrm{read}^{-1}(\mathrm{read}(\omega))\).
On the contrary, if \(\omega'\) ​​satisfies \(\mathrm{read}(\omega')=\mathrm{read}(\omega)\), and the bijection \(g\) of \(\omega\) and \(\omega'\) is exchanged (the other points are identical), then \(g\in\mathrm{Aut}_{\mathrm{read}}(\Omega)\) and \(g\omega=\omega'\) are defined.
Certification completed.
\end{proof}

\begin{proposition}[\texorpdfstring{$\mathsf{NULL}$}{NULL}’s three-decomposition and injectivization principle]\label{prop:xi-null-three-way}
Suppose \(\pi:\Omega\to Y\) ​​is a readout (many-to-one compression is allowed), and "checkable single-value readout" is used as the defining premise of the type layer according to the convention of "address before value" (see section \ref{sec:spg}).
If not provided enable extended read
\(
\pi^{\mathrm e}=(\pi,r):\Omega\to Y\times R
\)
Becomes the injective recording axis \(r\), then the readout degenerates to \(\mathsf{NULL}\) at the type level.
And the \(\mathsf{NULL}\) ​​can be decomposed into three types of sources according to the status of the record axis \(r\):
\begin{enumerate}
  \item ​​\textbf{Semantic null value:} The record axis is empty (the address is not given), so the readout is not defined;
  \item \textbf{Protocol null value:} The record axis is prohibited by the interface protocol/type system (for example, the radial mode of \(|u|\neq 1\) ​​is rejected under the unitary slice scope), so the object cannot be typed;
  \item \textbf{Collision null:} The record axis exists but the budget is not enough to achieve injectiveization; under the interval readout/CRT external protocol (defined \ref{def:pom-order-elimination-rewrite}), if you want to promote the folded many-to-one compression to be approximately reversible in the sense of failure probability \(\le\varepsilon\), the modular product budget of prime-register must satisfy
  \[
  P_m>2B_m(\varepsilon),
  \qquad
  B_m(\varepsilon):=\exp(m\gamma_m(\varepsilon)),
  \]
  Where \(\gamma_m(\varepsilon)\) ​​is the tail position of the fiber index (definition \ref{def:pom-gamma-quantile}).
  When \(P_m>2B_m(\varepsilon)\), the failure probability of "label out-of-bounds leading to non-unique reconstruction" can be reduced to \(\le\varepsilon\); and this threshold is sharp in the sense of this protocol (corollary \ref{cor:pom-prime-register-quantile-budget}).
\end{enumerate}
\end{proposition}

\begin{remark}[Register minimality and “end patch” incompleteness]\label{rem:xi-register-minimality-no-tail-patch}
The repair of collision nulls is not accomplished by a local patch of "appending a few bits at the end": in resolution reduction, the difference between naive truncation and fold-aware restriction can be as large as a full-length flip under the worst input (Theorem \ref{thm:fold-gauge-anomaly-max}), and any patch that "only repairs the terminal constant bits" is incomplete (corollary \ref{cor:fold-tail-patch-incomplete}).
Therefore, once a certain type of information really needs to be externalized as side information to ensure reversibility or verifiability, its record axis must be \emph{global non-local} in structure: it can only fall on the residue/inverse limit axis carried by the prime-register (and the finite state memory consistent with folding), and cannot be used as an additional continuous coordinate register decoupled from the phase \(\theta\). In this paper \(\mathbf{PW}\) Occurs within a closure interface.
\end{remark}

\begin{theorem}[\texorpdfstring{$\mathbf{PW}$}{PW} Continuous Hair Principle of Closures (Type—Register Minimality)]\label{thm:xi-pw-no-continuous-hair}
Under the visible domain constraints (\(X_{\mathrm{vis}}=\mathsf{Phase}\)) defining \ref{def:xi-visible-read-null} and the \(\mathbf{PW}\) regularized closure audit interface used In this paper, the following two structural assertions are established:
\begin{enumerate}
  \item ​​\textbf{(Radial non-absorbable)}
  If the reference/verification of an object structurally requires the radial data of \(|u|\neq 1\) ​​(equivalently requires the radial-phase double parameter structure) to keep its pole/main scale unchanged, then under the pure phase system, the object has no address within the closure, and its readout degenerates to \(\mathsf{NULL}\) according to the protocol;
  And without introducing additional recording axes, no constant layer strictification can eliminate this radial dependence while keeping the pole/primary scale unchanged.
  \item ​​\textbf{(External non-locality and minimality)}
  If the goal is to reduce \(\mathsf{NULL}\) ​​and allow external side information to achieve injectiveization of extended readout, then the external axis must be a global non-local residue/inverse limit axis (prime-register and its consistency constraints) in structure, and cannot be implemented by an additional continuous coordinate register decoupled from the phase \(\theta\).
\end{enumerate}
\end{theorem}

\begin{proof}
\textbf{(1)} is defined by \ref{def:xi-visible-read-null}. Under the visible domain \(X_{\mathrm{vis}}=\mathsf{Phase}\), any object that requires \(|u|\neq 1\) does not meet the visible domain filtering conditions, so its readout degenerates to \(\mathsf{NULL}\) by definition.
On the other hand, the local reversible readout of the radial-phase state \((\theta,\varrho)\) ​​cannot be continuously compressed into a single circumferential register (Theorem \ref{thm:radial-dimension-obstruction}). Therefore, under the system that rejects the external recording axis, "elimination of radial degrees of freedom with a constant layer" cannot be established while keeping the pole/main scale unchanged; its semantic consequence is that it is not addressable (\(\mathsf{NULL}\)).
\par
\textbf{(2)} According to the proposition \ref{prop:xi-null-three-way}, if there is folding/many-to-one compression in the readout, the record axis that makes the extended readout become injective must be external, otherwise it will degenerate to \(\mathsf{NULL}\) ​​at the type level; and the prime-register budget gives an auditable threshold to achieve injectiveization.
Furthermore, register minimality and "end-patch" incompleteness indicate that any record axis that is sufficient to repair collision nulls must be globally non-local in structure, and can only fall on the residue/inverse limit axis (note \ref{rem:xi-register-minimality-no-tail-patch}).
In summary, two assertions are made. Certification completed.
\end{proof}

\paragraph{\(\mathbf{PW}\) ​​The “visible” form of closures and interface parameters}
On the three-generator fragment \((\mathrm{PROJ}_u,E_m,\mathrm{LIFT}_G)\), any projected word can be reduced to the interface regular form
\(\mathrm{LIFT}_G\circ\mathrm{PROJ}_u\circ E_m\), and completely wrap up the non-commutative residual into the abnormal signature \(\mathsf{Anom}_G(K;\theta)\) (proposition \ref{prop:pom-three-gen-interface-normal-form} and definition \ref{def:pom-anom-signature}).
When this normal form is connected to the \(\Xi\)-audit interface in this chapter, the continuous parameter axis that can be reviewed is the phase of the unitary slice \(\theta\) (i.e. \(u=e^{i\theta}\)), rather than the general radial-phase decomposition \(u=re^{i\theta}\).
Therefore, the "addressable \(\Xi\) ​​zero-point certificate" can only be presented in the \(\theta\)-scope in the closure of this paper.

\begin{theorem}[\(\mathbf{PW}\) ​​Type safety: non-addressability of offset points (unitary slice diameter)]\label{thm:xi-offset-null-type-safety}
Pin \(L>1\) ​​and consider \(\Xi_{\Delta,L}(s)=\Delta(L^{-s},L^{2s-1})\) (as defined in \S\ref{subsec:xi-from-self-dual}).
In the "three-generator regularizer + exception signature interface + unitary slice zero-point audit" closure used In this paper, if \(s_0\in\CC\) ​​satisfies
\(\Re(s_0)\neq \tfrac12\)，
Then \(s_0\) ​​is not addressable under this closure; equivalently, the corresponding readout is a null value (denoted as \(\mathsf{NULL}\)).
In particular, any \(\Xi_{\Delta,L}\) ​​zero point \(s\) that is addressable and checkable within this closure must satisfy \(\Re(s)=\tfrac12\).
\end{theorem}

\begin{proof}
Got by \(u_L(s)=L^{2s-1}\)
\[
|u_L(s_0)|=L^{2\Re(s_0)-1}\neq 1,
\]
So \(u_L(s_0)\notin\mathsf{Phase}\).
\par
On the other hand, the zero-point audit interface in this chapter is parameterized with \(u=e^{i\theta}\) ​​(\(\theta\in\RR\)) on the unitary slice (proposition \ref{prop:xi-unitary-slice-spectral-crossing}), and accordingly the zero-point detection is converted into an auditable spectrum crossing criterion (note \ref{rem:xi-unitary-slice-operator-audit}).
On the projected word side, three-generator regularization converges any candidate compiled word into the interface canonical form and compresses the residual into a \(\theta\)-parameterized anomaly signature (the proposition \ref{prop:pom-three-gen-interface-normal-form} and the definition \ref{def:pom-anom-signature}), so that the visible continuous parameter axis of the "verifiable certificate" is aligned with the \(\theta\)-diameter of the unitary slice mentioned above.
\par
Therefore, within the closure of this paper, if a certain point \(s_0\) ​​can be addressed as an \(\Xi\) audit object, then its parameter \(u_L(s_0)\) must fall in the visible parameter domain \(\mathsf{Phase}\); but the above formula shows \(u_L(s_0)\notin\mathsf{Phase}\), a contradiction.
Therefore, \(s_0\) ​​has no address under this closure. According to the null value semantics in Section \ref{sec:recursive-addressing}, its readout is not defined, which is the null value \(\mathsf{NULL}\).
\par
Finally, "the addressable zero point must be on the critical line" is established by the inverse negation: if \(s\) ​​is addressable as the zero point within the closure, it is impossible to satisfy \(\Re(s)\neq\tfrac12\), so \(\Re(s)=\tfrac12\). Certification completed.
\end{proof}

\begin{corollary}[Critical linearity of visible zero point (internal RH of \(\Xi\))]\label{cor:xi-visible-zeros-on-critical-line}
Under the closure setting of theorem \ref{thm:xi-offset-null-type-safety}, all addressable and checkable \(\Xi_{\Delta,L}\) zeros in \(\mathbf{PW}\) ​​satisfy \(\Re(s)=\tfrac12\).
\end{corollary}

\begin{remark}[The difference between soundness and completeness]\label{rem:xi-soundness-vs-completeness}
Theorem \ref{thm:xi-offset-null-type-safety} is a type-safe soundness direction: \emph{addressable/auditable} \(\Rightarrow\) ​​\emph{critical linearity}.
And proposition
\[
\Xi_{\Delta,L}(s)=0\ \Longrightarrow\ s\ \text{at}@@TK1 3@@\mathbf{PW}\ \text{medium addressable (not}\ \mathsf{NULL}\text{)}
\]
is the completeness direction, and its content is not equivalent to "all visible concepts are non-null" (the latter is a defining statement in the semantics of section \ref{sec:recursive-addressing}).
When the internal scope needs to be upgraded to \(\Xi\)-RH, the gap in completeness must be filled: an auditable certificate condition is given to ensure that the zero-point event is addressable under the closure (without triggering \(\mathsf{NULL}\)).
\par
Under the unit circle phase semantics, the above two can also be understood as the distinction between two "\(\mathsf{NULL}\)---conservation" systems.
\emph{Pure phase regime} takes \(X_{\mathrm{vis}}=\mathsf{Phase}\) ​​and replaces it with \(\mathsf{NULL}\) to maintain type safety: any object requiring \(|u|\neq 1\) radial degrees of freedom is considered unaddressable and thus rejected within the visible semantics (corresponding to soundness in this section).
If the goal is to reduce the occurrence of \(\mathsf{NULL}\), the squashed degrees of freedom must be externalized according to the "expansion preservation" principle (proposition \ref{prop:unit-circle-phase-extension-preserving}):
Write \(u=re^{i\theta}\) ​​and get
\[
\operatorname{proj}_{\mathsf{Phase}}(u)=\frac{u}{|u|}=e^{i\theta}\in\mathsf{Phase},
\qquad
\varrho(u)=\log|u|\in\RR,
\qquad
\operatorname{proj}^{\mathrm{e}}(u):=\bigl(\operatorname{proj}_{\mathsf{Phase}}(u),\varrho(u)\bigr).
\]
Then \(\operatorname{proj}^{\mathrm{e}}\) ​​is injective on \(\CC^\times\), and the unitary slice is only the section \(\varrho(u)=0\);
At this time, \(|u|\neq 1\) ​​no longer necessarily triggers \(\mathsf{NULL}\), but behaves as a non-degraded extended readout of "phase readout + radial recording".
When the recording axis is allowed to grow without an upper bound, the external degrees of freedom described above can be carried by additional orthogonal modes of the separable Hilbert space in the limiting representation of a "compact phase tower" (Proposition \ref{prop:phase-tower-principal-u1}).
Therefore, the completeness gap required for RH can be located as: \emph{Is there a zero-point event that must be introduced to explain the radial channel (or other additional visible degrees of freedom, such as the minimum realized dimension exceedance)}; if it exists, it must read as \(\mathsf{NULL}\) ​​in the pure phase regime, while in the extended regime the additional recording axis is required to explicitly carry its radial residual.
\end{remark}

\begin{remark}[Intrinsic factor hidden component and \texorpdfstring{$\mathsf{NULL}$}{NULL}: Compilability is typing]\label{rem:xi-inner-factor-auditability-null}
Under the Hardy/Smirnov conditional aperture of diskization completion \(F_{\zeta}\), the boundary mode length \(|F_{\zeta}(e^{\mathrm{i}\theta})|\) only determines the external factor \(O\), while the zero point data in the disk is carried by the Blaschke internal factor \(B\) (\S\ref{subsec:unit-disk-inner-outer}).
In particular, the offset from the critical line \(\delta=\tfrac12-\beta\) ​​is strictly implemented in the disk as a horocycle/Busemann coordinate with the endpoint \(-1\) as the base point
\(
\delta=B^{-1}(w_{\rho})=\frac{1-|w_{\rho}|^{2}}{|1+w_{\rho}|^{2}}
\)
(Corollary \ref{cor:app-offcritical-horocycle-busemann}), and therefore belongs to the hidden component \emph{structurally invisible} of unitary slice pure phase readout.
\par
In this paper's \(\mathbf{PW}\) ​​closure audit, "whether hidden components are acceptable" is a type issue rather than a numerical issue:
If the external record axis is rejected, then according to the theorem \ref{thm:xi-pw-no-continuous-hair} and proposition \ref{prop:xi-null-three-way}, this hidden degree of freedom will trigger collision/protocol null value and degenerate in the form of \(\mathsf{NULL}\);
If it is required to remain verifiable within a closed loop, the intrinsic factor data must be compiled into auditable output (such as the resolution threshold and angular atoms of the defect ledger \(D_{\zeta}(\varrho)\), or discrete spectral fingerprints such as the Hankel rank/rigidity index of the implementation layer), making it enter the limited semantic output domain of the "signature \(\times\) spectrum" (definition \ref{def:pom-sem-functor}).
Therefore, in this system, "offline information is not visible" can be strictly simplified as: \emph{If it cannot be certified, it cannot be typed (read as \(\mathsf{NULL}\))}}.
\end{remark}

\begin{remark}[Auditable Teardown of Completeness Gap: Cross-Resolution Patch and Time Layer Right Decoding]\label{rem:xi-completeness-auditable-decomposition}
completeness can be expressed as a global assertion that "the zero-point event does not trigger \(\mathsf{NULL}\)".
In the projection gate semantics of this paper, this assertion can be further decomposed into several auditable sub-conditions, so that the source of \(\mathsf{NULL}\) ​​can be isolated layer by layer.
\begin{enumerate}
  \item ​​\textbf{Cross-resolution: Incomplete finite tail patch (cannot be partially eliminated). }
  Microstate direct truncation \(r_{m+1\to m}\) ​​is generally incompatible with fold-aware restriction \(\widetilde r_{m+1\to m}\), and specification difference \(G_m\) (definition \ref{def:fold-gauge-anomaly}) quantifies its minimum correction size;
  And any "fix only the terminal constant bits" patch is impossible to reproduce \(\widetilde r_{m+1\to m}\) ​​on all inputs (corollary \ref{cor:fold-tail-patch-incomplete}; and see note \ref{rem:fold-tail-patch-null-nonlocal}).
  Therefore, if the protocol denies both non-local access and orthogonal record axes, some input families will not be able to obtain fold-consistent addresses, and their readouts will degenerate to \(\mathsf{NULL}\) ​​according to null value semantics.

  \item ​​\textbf{Time layer: right reversibility bisection law and synchronization budget. }
  After promoting \(\Fold_m\) ​​along the sliding window to the sofic factor \(\Phi_m\), the statistical type \(\mathsf{NULL}\) of the time layer can be characterized as a deterministic filtering state non-single point (\(|S_t|>1\); note \ref{rem:Ym-stat-null-unsync}).
  Its behavior satisfies the dichotomy of right reversibility: either it is fully synchronized after a unified finite delay (\(Y_m^{\mathrm{amb}}=\varnothing\), corollary \ref{cor:Ym-time-layer-dichotomy}; the strong form is given by the theorem \ref{thm:Ym-ambiguity-shell-dag} when \(m\ge 3\) is in the Zeckendorf fold), or long-term ambiguity can only reside on strictly low-entropy shells, with exponential synchronization tail bounds under the maximum entropy measure \(\nu_m(T_m>n)\le C_m e^{-\varepsilon_m n}\) (Theorem \ref{thm:Ym-sync-tail}, where \(\varepsilon_m=\log 2-h_{\mathrm{top}}(Y_m^{\mathrm{amb}})>0\)).

  \item ​​\textbf{Constant-level sealing: period merging and online risk table. }
  The main pole residue ratio \(\eta_m\in(0,1]\) of cover and intrinsic \(\zeta\) ​​(proposition \ref{prop:Ym-zeta-residue-ratio}) and the period compression defect \(\delta_m\in[0,\log|\mathcal{S}_m|]\) (corollary \ref{cor:Ym-periodic-compression-ratio}) quantify the constant-level information gap of "merging but not changing the exponential scale";
  The single-point state proportion limit \(\eta_m^{\mathrm{obs}}\in(0,1]\) ​​(proposition \ref{prop:eta_m_hat_ergodic_limit}) gives an online observable risk table without a determinant.
\end{enumerate}
\par
Therefore, the "no \(\mathsf{NULL}\)" goal of completeness can be replaced by a family of auditable criteria:
Once the time layer satisfies finite delay synchronization (or the synchronization tail budget is met) and constant-level merging is controlled (e.g. \(\eta_m^{\mathrm{obs}}\approx 1\) ​​and \(\eta_m\approx 1\)), if systematic \(\mathsf{NULL}\) still occurs, its source should be prioritized to higher-level mechanisms such as geometric out-of-domain (radial degrees of freedom of \(|u|\neq 1\)) or dimensionality thresholds (minimum realization over-bounds), rather than time-layer ambiguity or insufficient cross-resolution finite tail patches.
\end{remark}

\begin{theorem}[NULL’s orthogonal trichotomy and irreplaceability of resource threshold]\label{thm:xi-null-orthogonal-trichotomy-resource-nonsubstitutable}
Under the unitary slice visible domain constraints of this paper (definition \ref{def:xi-visible-read-null}, take \(X_{\mathrm{vis}}=\mathsf{Phase}\)), \(\mathbf{PW}\) closure type safety (theorem \ref{thm:xi-offset-null-type-safety}) and injectivity principle (proposition \ref{prop:xi-null-three-way}),
The structural sources of global readout failure (\(\mathsf{NULL}\)) can be organized into three types of witnesses that are not equivalent to each other and are \emph{orthogonal} in an auditable sense:
\begin{enumerate}
  \item ​​\textbf{(closure/out-of-scope defect: reference or type level rejection)}
  The address is not given, or the interpretation of the object requires degrees of freedom beyond \(X_{\mathrm{vis}}\) ​​(typically: the radial mode of \(|u|\neq 1\)), such that there is no address within the closure and is read by protocol as \(\mathsf{NULL}\) (definition \ref{def:xi-visible-read-null}; theorem \ref{thm:xi-pw-no-continuous-hair}).
  \item ​​\textbf{(Adhesive kink: second-order compatibility barrier)}
  Even if a local certificate exists slice-by-slice on the coverage, it may not be globalized due to a failure of the second-order consistency constraint; its failure can be witnessed by a verifiable Čech \(2\)-cocycle obstacle class \([\alpha]\in H^2(\mathrm{Addr},A)\) (definition \ref{def:prefix-site-h2-gluing-obstruction}), and it follows that the address fiber is empty and read as \(\mathsf{NULL}\) (proposition \ref{prop:null-as-h2-obstruction}).
  \item ​​\textbf{(Insufficient hard resources: threshold budget violation)}
  In order to upgrade many-to-one compressed readout to single-value verifiable readout, two types of irreplaceable hard threshold resources must be met at the same time:
  \begin{itemize}
    \item \emph{Collision threshold (prime-register budget):} If the modular product budget \(P_m\) ​​is not enough to exceed the quantile threshold \(2B_m(\varepsilon)\), then the collision null value cannot be eliminated under a given confidence scope (Proposition \ref{prop:xi-null-three-way}; and see the worst-case lower bound of theorem \ref{thm:address-defect-law}).  
    \item \emph{Endpoint threshold (accuracy/resolution budget):} If the endpoint layer resolution is not sufficient to carry the secondary compression gap of the off-slice signal, the offline information cannot be single-valued in the verification semantics and must appear in the form of \(\mathsf{NULL}\) or "the mode axis must be extended" (theorem \ref{thm:xi-offcritical-endpoint-resolution-lower-bound}; see also theorem \ref{thm:xi-radial-information-projection-lower-bound} The compactification reads out a lower bound on bit complexity).
  \end{itemize}
\end{enumerate}
And the above three types of witnesses are not interchangeable within the system: \emph {increase the budget of one axis} cannot replace \emph {fix the defects of another axis} (for example, the modulus increase of prime-register cannot offset the endpoint resolution threshold, and vice versa; for the same reason, any budget increase cannot eliminate the bonding obstacle of \([\alpha]\neq 0\)).
This irreplaceability is the auditable meaning of "information conservation/non-leakage" In this paper: the sources of failure are each closed under orthogonal decomposition, and each carries an independent threshold or invariant certificate.
\end{theorem}

\begin{proof}
Class (1) is a direct consequence of visible domain filtering and \(\mathbf{PW}\) ​​type safety (definition \ref{def:xi-visible-read-null}; theorem \ref{thm:xi-offset-null-type-safety} and theorem \ref{thm:xi-pw-no-continuous-hair}).
Class (2) is given by the \(H^2\) ​​adhesive barrier scope at the prefix site: if \([\alpha]\neq 0\), then there is no globalizable local section, so the fiber is empty and read as \(\mathsf{NULL}\) (Proposition \ref{prop:null-as-h2-obstruction}).
Class (3) is given by the collision threshold of the injectivity principle (Proposition \ref{prop:xi-null-three-way}) and the endpoint quadratic compression/resolution barrier (Theorem \ref{thm:xi-offcritical-endpoint-resolution-lower-bound}; Theorem \ref{thm:xi-radial-information-projection-lower-bound}) respectively.
\par
The irreplaceability comes from the fact that the "support axes" of these certificates do not overlap with each other: the collision threshold only constrains the modulus budget of the residue/inverse limit axis, while the endpoint threshold only constrains the resolution threshold induced by the endpoint precision potential; the two correspond to different orthogonal record axes in the type semantics, so there is no legal compilation path to use one type of budget to cover the gap of another type.
In the same way, \([\alpha]\neq 0\) ​​is an invariant certificate of second-order consistency and cannot be eliminated by any simple bit budget increase.
Certification completed.
\end{proof}

\input{sections/body/zeta_finite_part/xi/subsubsec__xi-null-z2-cech-transfer-doublecover-spectral-splitting}

\input{sections/body/zeta_finite_part/xi/cor__xi-us-offslice-dichotomy-null-decidability}

\paragraph{Tokens and Semantic Gates: Folding Fibers, Maximum Entropy Boost and Invisible Capacity}\label{par:xi-null-notation-semantic-gate}
The following notation is used to unify "addressable/null" and window-folded fiber information gaps on unitary slices into the same ledger language (see section \ref{subsec:pom-information-ledger} ​​and definition \ref{def:pom-kappa}).
Fixed resolution \(m\ge 1\), taking microscopic phase space \(\Omega_m:=\{0,1\}^m\) and stable type set \(X_m\).
Let the folding projection be \(\Fold_m:\Omega_m\to X_m\) ​​and the fiber multiplicity be
\[
d_m(x):=\card{\Fold_m^{-1}(x)}\qquad(x\in X_m).
\]
For any macroscopic edge distribution \(\pi\in\Delta(X_m)\), define its \textbf{fiber uniform lift} (maximum entropy lift)
\begin{equation}\label{eq:xi-null-maxent-lift}
\pi^{\mathrm{e}}(a):=\frac{\pi(\Fold_m(a))}{d_m(\Fold_m(a))},\qquad a\in\Omega_m,
\end{equation}
Thus \((\Fold_m)_*\pi^{\mathrm{e}}=\pi\); it is the same as \(\widetilde\pi\) where the text defines \ref{def:pom-maxent-lift}.
And define \textbf{thermodynamic defect/invisible capacity}
\[
\kappa_m(\pi):=\mathbb{E}_{\pi}\bigl[\log d_m(X)\bigr]
\]
(definition \ref{def:pom-kappa}).
In the unitary slice closure in this section, the "addressable" continuous field is fixed to \(\mathsf{Phase}\), and is additionally constrained by the \(\mathbf{PW}\) type closure (definition \ref{def:xi-visible-read-null});
Therefore, if external additional recording axes are rejected, the internal label of the folded fiber has no address in the visible semantics, and its "readout" is treated equivalently as \(\mathsf{NULL}\) ​​(unaddressable) at the semantic level.

\begin{theorem}[Micro entropy window and endpoint implementation under fixed macro \texorpdfstring{$\pi$}{pi}]\label{thm:xi-null-fiber-entropy-window}
Fixed resolution \(m\ge 1\) ​​and macro edge \(\pi\in\Delta(X_m)\), and made \(\pi^{\mathrm{e}}\in\Delta(\Omega_m)\) uniformly enhance its fibers \eqref{eq:xi-null-maxent-lift}.
Then for any lift \(\mu\in\Delta(\Omega_m)\) ​​that satisfies \((\Fold_m)_*\mu=\pi\), a strict ledger identity is established
\[
\boxed{\;
H(\mu)=H(\pi)+\kappa_m(\pi)-D_{\mathrm{KL}}(\mu\|\pi^{\mathrm{e}})\; }.
\]
and:
\begin{enumerate}
  \item \(\pi^{\mathrm{e}}\) ​​\textbf {uniquely} achieves maximum entropy among all lifts, and
  \[
  H(\pi^{\mathrm{e}})=H(\pi)+\kappa_m(\pi).
  \]
  \item There is a “fiber point mass lift” \(\mu^\star\) ​​such that \(H(\mu^\star)=H(\pi)\) and
  \[
  D_{\mathrm{KL}}(\mu^\star\|\pi^{\mathrm{e}})=\kappa_m(\pi).
  \]
  Therefore, under a fixed \(\pi\), the microscopic entropy satisfies \textbf{exact window width}
  \[
  H(\mu)\in\bigl[H(\pi),\ H(\pi)+\kappa_m(\pi)\bigr],
  \qquad
  \mathrm{width}=\kappa_m(\pi).
  \]
\end{enumerate}
\end{theorem}

\begin{proof}
The ledger identity is the “fixed \(\pi\)” version of theorem \ref{thm:pom-kl-ledger} (where \(\mu^e\) is denoted as \(\pi^{\mathrm{e}}\)).
The maximum entropy endpoint and uniqueness are given by the theorem \ref{thm:pom-maxent-lift}.
\par
Point mass endpoint: Pick any section \(s:X_m\to\Omega_m\) ​​that satisfies \(\Fold_m(s(x))=x\) (just pick any point for each fiber).
definition
\(
\mu^\star(s(x)):=\pi(x)
\)
And for the remaining \(a\notin s(X_m)\), take \(\mu^\star(a)=0\).
Then \((\Fold_m)_*\mu^\star=\pi\), and
\[
H(\mu^\star)=-\sum_{x\in X_m}\pi(x)\log\pi(x)=H(\pi).
\]
And from \(\pi^{\mathrm{e}}(s(x))=\pi(x)/d_m(x)\), we get
\[
D_{\mathrm{KL}}(\mu^\star\|\pi^{\mathrm{e}})
=\sum_{x\in X_m}\pi(x)\log\frac{\pi(x)}{\pi(x)/d_m(x)}
=\sum_{x\in X_m}\pi(x)\log d_m(x)
=\kappa_m(\pi).
\]
Certification completed.
\end{proof}

\begin{theorem}[Address—Defect Law: Side Information Budget Required for Injection]\label{thm:address-defect-law}
Fixed resolution \(m\ge 1\) ​​with folded readout \(\Fold_m:\Omega_m\to X_m\), and made the recording axis \(r:\Omega_m\to R\).
\begin{enumerate}
  \item ​​\textbf{(worst case: fiber pointwise lower bound)}
  If the extended readout \((\Fold_m,r):\Omega_m\to X_m\times R\) ​​is injective, then for each \(x\in X_m\)
  \[
  \card{r(\Fold_m^{-1}(x))}\ \ge\ d_m(x),
  \]
  Thus \(\card{R}\ge \max_{x\in X_m} d_m(x)\).
  In particular, as long as there is \(x\) ​​and \(d_m(x)>1\), the stable type \(x\) alone cannot form a single-value verifiable "complete address"; if the external record axis is rejected, the collision corresponds to the collision type \(\mathsf{NULL}\) under the semantics of this paper (definition \ref{def:xi-visible-read-null} and proposition \ref{prop:xi-null-three-way}).

  \item ​​\textbf{(Thermodynamics: average bit budget lower bound)}
  Given any macro edge \(\pi\in\Delta(X_m)\), let \(A\sim\pi^{\mathrm{e}}\in\Delta(\Omega_m)\) be uniformly lifted by its fibers \eqref{eq:xi-null-maxent-lift},
  And order \(X:=\Fold_m(A)\sim\pi\).
  If \((\Fold_m,r)\) ​​is injective, then
  \[
  H\bigl(r(A)\mid X\bigr)\ \ge\ H(A\mid X)\ =\ \kappa_m(\pi),
  \]
  Therefore, the recording axis needs to carry at least \(\kappa_m(\pi)/\log 2\) ​​bits of side information in an average sense.
\end{enumerate}
\end{theorem}

\begin{proof}
\textbf{(1)} Fixed \(x\in X_m\). If \(a\neq a'\) and \(\Fold_m(a)=\Fold_m(a')=x\), the injectivity of \((\Fold_m,r)\) must be \(r(a)\neq r(a')\), so \(r\) is injective on fiber \(\Fold_m^{-1}(x)\), and then \(\card{r(\Fold_m^{-1}(x))}\ge \card{\Fold_m^{-1}(x)}=d_m(x)\). Take the maximum value to get \(\card{R}\ge \max_x d_m(x)\).
\par
\textbf{(2)} Since \(\pi^{\mathrm{e}}\) ​​is uniform on each fiber, for each \(x\)
\(
H(A\mid X=x)=\log d_m(x)
\), thus
\(
H(A\mid X)=\mathbb{E}_\pi[\log d_m(X)]=\kappa_m(\pi)
\)。
On the other hand, \((\Fold_m,r)\) ​​is injective meaning there is a deterministic decoder \(D\) that makes \(A=D(X,r(A))\), so \(H(A\mid X,r(A))=0\).
then
\[
H(A\mid X)=I(A;r(A)\mid X)\le H(r(A)\mid X),
\]
Obtains the lower bound shown; divide by \(\log 2\) ​​to convert bit units. Certification completed.
\end{proof}

\begin{remark}[Semantic explanation: Three types of sources of \texorpdfstring{$\mathsf{NULL}$}{NULL} under unitary slice]\label{rem:xi-null-three-kinds}
In unitary slice audit semantics, if external orthogonal record axes are rejected, "unaddressable/null values" can be decomposed into three types of mutually orthogonal structural sources:
\begin{itemize}
  \item \textbf{Geometry type \(\mathsf{NULL}\).} The parameter domain constraint is \(\mathsf{Phase}\): Any object that requires \(|u|\neq 1\)’s radial degrees of freedom has no address within the closure, and by definition \ref{def:xi-visible-read-null} is read as \(\mathsf{NULL}\).
  \item \textbf{Branch type \(\mathsf{NULL}\).} Even if it is limited to \(|u|=1\), if the half-angle layer \(\mathfrak r\) is further pressed to the only angle layer \(u=\mathfrak r^2\), the discrete branch that double covers \(\mathfrak r\mapsto-\mathfrak r\) will be erased; the branch information is not visible to the \(u\)-layer (see appendix \S\ref{subsec:unit-disk-witt-chebyshev} and proposition \ref{prop:completion-subring-unique-angle}).
  \item \textbf{Statistical \(\mathsf{NULL}\).} When fixing the macro edge \(\pi\), lift is not unique; if there is no external fiber shape record, the micro difference has no address in the visible semantics.
  Theorem \ref{thm:xi-null-fiber-entropy-window} gives its invisible window width \(\kappa_m(\pi)\), and accurately carries the "fiber shape residual" of any specific lift as \(D_{\mathrm{KL}}(\mu\|\pi^{\mathrm{e}})\).
\end{itemize}
\end{remark}

\begin{corollary}[Fiber residual decomposition and “statistical \texorpdfstring{$\mathsf{NULL}$}{NULL} radius”]\label{cor:xi-null-statistical-radius}
Follow the notation of theorem \ref{thm:xi-null-fiber-entropy-window} and write fiber conditional distribution \(\mu_x(\cdot):=\mu(\cdot\mid x)\) ​​(\(x\in X_m\)) and fiber uniform distribution \(u_x(a):=1/d_m(x)\) (\(a\in\Fold_m^{-1}(x)\)).
Then there is fiber orthogonal decomposition
\[
\boxed{\;
D_{\mathrm{KL}}(\mu\|\pi^{\mathrm{e}})
=\sum_{x\in X_m}\pi(x)\,D_{\mathrm{KL}}(\mu_x\|u_x)\; }.
\]
Thus there is an explicit upper bound on all lift
\[
0\le D_{\mathrm{KL}}(\mu\|\pi^{\mathrm{e}})\le \kappa_m(\pi),
\qquad
D_{\mathrm{TV}}(\mu,\pi^{\mathrm{e}})\le \sqrt{\frac{\kappa_m(\pi)}{2}}.
\]
And for any bounded observable \(f:\Omega_m\to\RR\), the systematic deviation of all lifts compatible with \(\pi\) on \(f\) relative to \(\pi^{\mathrm{e}}\) satisfies
\[
\bigl|\mathbb{E}_\mu[f]-\mathbb{E}_{\pi^{\mathrm{e}}}[f]\bigr|
\le 2\|f\|_\infty\sqrt{\frac{\kappa_m(\pi)}{2}}.
\]
\end{corollary}

\begin{proof}
The first formula is a direct restatement of Lemma \ref{lem:pom-fiberwise-kl} (denoting \(\mu^e\) ​​as \(\pi^{\mathrm{e}}\)).
The remaining inequalities are established by theorem \ref{thm:defect-limited-audit} ​​(and Pinsker’s inequality). Certification completed.
\end{proof}

\begin{corollary}[Microscopic sliding window: KL interval law under the same macroscopic view]\label{cor:xi-null-micro-kl-window}
Let \(u_{\Omega_m}\) ​​be the uniform distribution on \(\Omega_m\), and let
\(
w_m:=(\Fold_m)_*u_{\Omega_m}
\)
Push it forward (definition \ref{def:pom-wm-baseline}).
Pin any \(\pi\in\Delta(X_m)\).
Then for all lift \(\mu\in\Delta(\Omega_m)\) that satisfies \((\Fold_m)_*\mu=\pi\), explicit clamping is established
\[
\boxed{\;
D_{\mathrm{KL}}(\pi\|w_m)
\le D_{\mathrm{KL}}(\mu\|u_{\Omega_m})
\le D_{\mathrm{KL}}(\pi\|w_m)+\kappa_m(\pi)\; }.
\]
Equivalently,
\(
0\le D_{\mathrm{KL}}(\mu\|\pi^{\mathrm{e}})\le \kappa_m(\pi)
\)。
\end{corollary}

\begin{proof}
From the identity \(D_{\mathrm{KL}}(\mu\|u_{\Omega_m})=m\log 2-H(\mu)\) ​​and substituting the theorem \ref{thm:xi-null-fiber-entropy-window} we get
\[
D_{\mathrm{KL}}(\mu\|u_{\Omega_m})
=m\log 2-H(\pi)-\kappa_m(\pi)+D_{\mathrm{KL}}(\mu\|\pi^{\mathrm{e}}).
\]
Then from the proposition \ref{prop:pom-kappa-kl-decomposition}
\(
\kappa_m(\pi)=m\log 2-H(\pi)-D_{\mathrm{KL}}(\pi\|w_m)
\)
get decomposition
\(
D_{\mathrm{KL}}(\mu\|u_{\Omega_m})
=D_{\mathrm{KL}}(\pi\|w_m)+D_{\mathrm{KL}}(\mu\|\pi^{\mathrm{e}})
\)。
Finally, apply the Corollary to the bounds of \ref{cor:xi-null-statistical-radius} for \(D_{\mathrm{KL}}(\mu\|\pi^{\mathrm{e}})\). Certification completed.
\end{proof}

\begin{remark}[Endpoint compression, limited precision and type 2 \texorpdfstring{$\mathsf{NULL}$}{NULL}]\label{rem:xi-null-finite-precision}
Unit circle area phase gates compress non-compact scales into the boundary layer: for example the height parameter \(t\in\RR\) ​​is squashed to \(u_{\mathrm C}(t)=\frac{1+\mathrm{i}t}{1-\mathrm{i}t}\in\TT\setminus\{-1\}\) on the Cayley axis, and \(t\to\pm\infty\) merges into the endpoint \(u_{\mathrm C}\to-1\) (Appendix \ref{app:unit-circle-phase-gate}).
At a strictly mathematical level, this gate is reversible and does not produce null values; however, in limited-precision implementation, there is strong error amplification in the back-projection of the endpoint neighborhood:
\[
\frac{\dd x}{\dd u}=\pi(1+x^2)
\qquad\text{(Formula \eqref{eq:app-precision-amplification})}.
\]
Therefore, when the accuracy budget is insufficient to distinguish different preimages in the endpoint neighborhood, the "readout" will no longer have verifiable single-valuedness: it will either degenerate into multi-valued/unverifiable, or be equivalently treated as unaddressable in the semantics of this paper (readout is undefined, i.e. \(\mathsf{NULL}\)).
In "extended save" language, this corresponds to the fact that the external recording axis (or precision register) is not sufficient to enable the extended observation to be injective, so that the address must be missing.
\end{remark}

\subsubsection{相位前缀过滤与二次分辨率门槛：端点压缩诱导的强制碰撞律}\label{subsubsec:xi-phase-prefix-resolution-law}
本段把“unitary slice / legal address / scan-to-clarify”的一参数读出语义具体化为：在可见轴仅允许相位紧化坐标的前提下，读出精细化由一个离散的\emph{分辨率深度}（前缀长度）度量，而非由额外连续坐标轴提供；任何离切片信息若未被外置到正交寄存器，则在对象层被迫塌缩为不可唯一读出（在本文语义中表现为碰撞型 \(\mathsf{NULL}\)，见定义 \ref{def:xi-visible-read-null} 与命题 \ref{prop:xi-null-three-way}）。
\par
关键结构事实在于：相位紧化把高高度区间压入端点邻域并产生二次压缩；当零点密度随高度以 \(T\log T\) 增长时，该二次压缩强制产生一个不可绕过的“二次分辨率屏障”。本段给出该屏障的可核验形式，并说明所谓“一个零点是另一个零点的地址”只能被严格化为“共享前缀所诱导的柱事件聚类”，而非点对点寻址。

\paragraph{相位盒与 dyadic 柱前缀}
取无参相位盒紧化门（附录 \ref{app:unit-circle-phase-gate}）
\[
u=\upsilon(t):=\frac{1}{\pi}\arctan(t)\in\Bigl(-\tfrac12,\tfrac12\Bigr),
\qquad
\widetilde u(t):=u+\tfrac12\in(0,1).
\]
对整数 \(m\ge 1\)，定义 \textbf{\(m\)-比特相位地址（合法读出前缀）}为 dyadic 柱单元
\[
I_m(t):=\Bigl[\frac{\lfloor 2^m\widetilde u(t)\rfloor}{2^m},\frac{\lfloor 2^m\widetilde u(t)\rfloor+1}{2^m}\Bigr)\subset(0,1),
\]
并记地址值
\(
A_m(t):=\lfloor 2^m\widetilde u(t)\rfloor\in\{0,1,\dots,2^m-1\}.
\)
在“地址先于取值”的语义下，\(m\) 即为\textbf{分辨率深度}：读出接口不引入新的连续自由度，只允许沿 \(m\uparrow\) 的柱过滤细化。

\paragraph{端点压缩下的地址碰撞强制律}
记 \(N_\zeta(T)\) 为 \(\zeta\) 在临界带 \(0<\Re(s)<1\) 内、满足 \(0<\Im(s)\le T\) 的非平凡零点计数（按重数计）。其经典主项满足
\[
N_\zeta(T)=\frac{T}{2\pi}\log\Bigl(\frac{T}{2\pi}\Bigr)-\frac{T}{2\pi}+O(\log T),
\qquad(T\to\infty).
\]

\begin{theorem}[端点压缩的地址碰撞下界：二次分辨率门槛]\label{thm:xi-phase-prefix-collision-lower-bound}
存在常数 \(T_0>0\) 与 \(c>0\)，使得对所有 \(T\ge T_0\) 与任意整数 \(m\ge 1\)，在高度带
\[
\Gamma(T):=\{\gamma>0:\ \exists\,\rho=\beta+\mathrm{i}\gamma\ \text{为非平凡零点}\ \text{且}\ T<\gamma\le 2T\}
\]
中，必存在某一相位地址值 \(a\in\{0,1,\dots,2^m-1\}\) 满足
\[
\#\{\gamma\in\Gamma(T):A_m(\gamma)=a\}\ \ge\ c\,\frac{T^2\log T}{2^m}.
\]
特别地，若希望高度带 \((T,2T]\) 上的相位地址映射 \(\gamma\mapsto A_m(\gamma)\) 不发生碰撞（单射），则必须满足
\[
2^m\ \gtrsim\ T^2\log T,
\qquad\text{等价地}\qquad
m\ \ge\ 2\log_2 T+\log_2\log T+O(1).
\]
\end{theorem}

\begin{proof}
首先，带 \([T,2T]\) 在 \(\upsilon\)-轴上的像长度满足
\[
\upsilon(2T)-\upsilon(T)=\frac{1}{\pi}\int_T^{2T}\frac{\dd t}{1+t^2}
\le \frac{1}{\pi}\int_T^{2T}\frac{\dd t}{t^2}
=\frac{1}{2\pi T}.
\]
因此深度为 \(m\) 的 dyadic 划分在该像带内能提供的不同地址单元个数至多为
\[
M_m(T)\le 2^m\bigl(\upsilon(2T)-\upsilon(T)\bigr)+2\ \le\ \frac{2^m}{2\pi T}+2.
\]
另一方面，由零点计数主项可得 \(N_\zeta(2T)-N_\zeta(T)\asymp T\log T\)（\(T\to\infty\)）。
将 \(\#\Gamma(T)=N_\zeta(2T)-N_\zeta(T)\) 个零点高度投入至多 \(M_m(T)\) 个地址盒，鸽巢原理给出存在某一盒占用数至少为
\[
\frac{\#\Gamma(T)}{M_m(T)}\ \gtrsim\ \frac{T\log T}{2^m/T}\ =\ \frac{T^2\log T}{2^m},
\]
即得第一式。单射必要条件由“最大占用数 \(\le 1\)”立得。
\end{proof}

\begin{corollary}[清晰度视界：固定分辨率深度的可分辨高度尺度]\label{cor:xi-phase-prefix-clarity-horizon}
固定 \(m\ge 1\)。若要求高度带 \((T,2T]\) 内相位地址无碰撞，则必有
\[
T^2\log T\ \lesssim\ 2^m.
\]
因此在忽略对数修正的主尺度下，可分辨高度视界满足 \(T\approx 2^{m/2}\)（并带 \(1/\sqrt{\log T}\) 级修正）。
\end{corollary}

\paragraph{“零点互为地址”的严格化：共享前缀而非点对点寻址}
对 \(u,v\in(0,1)\) 定义 dyadic 前缀共享长度
\[
L(u,v):=\max\{m\ge 0:\ \lfloor 2^m u\rfloor=\lfloor 2^m v\rfloor\}.
\]
令 \((\gamma_n)_{n\ge 1}\) 为非平凡零点的正虚部按递增排列（按重数计），并写 \(u_n:=\widetilde u(\gamma_n)\in(0,1)\)。

\begin{proposition}[共享前缀的二次阈值：相邻零点在端点层形成稳定簇]\label{prop:xi-phase-adjacent-prefix-sharing}
存在常数 \(T_1>0\) 与 \(C>0\)，使得对所有 \(T\ge T_1\)，在满足 \(T<\gamma_n\le 2T\) 的相邻零点对 \((\gamma_n,\gamma_{n+1})\) 中，至少有一半满足
\[
L(u_n,u_{n+1})\ \ge\ 2\log_2 T+\log_2\log T - C.
\]
\end{proposition}

\begin{proof}
考虑相位窗 \([\upsilon(T),\upsilon(2T)]\)。
其总长度由上证得 \(\upsilon(2T)-\upsilon(T)\ll 1/T\)。
而窗内相邻对的相位间隔之和恰为该总长度，故平均相位间隔满足
\[
\overline{\Delta u}(T)\ :=\ \frac{\upsilon(2T)-\upsilon(T)}{N_\zeta(2T)-N_\zeta(T)}\ \ll\ \frac{1/T}{T\log T}\ =\ \frac{1}{T^2\log T}.
\]
若超过一半的相邻对满足 \(|u_{n+1}-u_n|>2\overline{\Delta u}(T)\)，则总和将超过相位窗长度，矛盾。
因此至少一半相邻对满足 \(|u_{n+1}-u_n|\ll (T^2\log T)^{-1}\)。
对这样的 \(u\neq v\)，由 dyadic 前缀定义有 \(L(u,v)\ge \lfloor \log_2(1/|u-v|)\rfloor -1\)，代入即得结论。
\end{proof}

\begin{remark}[共享前缀 \(\neq\) 互为寻址：地址对象是柱事件而非点指针]\label{rem:xi-phase-prefix-not-pointer}
对固定深度 \(m\)，关系
\[
\gamma\sim_m\gamma'\quad:\Longleftrightarrow\quad A_m(\gamma)=A_m(\gamma')
\]
刻画的是“同一柱事件（同一 dyadic 地址盒）”，而非“点对点指针”。
定理 \ref{thm:xi-phase-prefix-collision-lower-bound} 表明：当 \(m<2\log_2T+\log_2\log T-O(1)\) 时，高度带 \((T,2T]\) 内该等价类在结构上必为高重数；因此任何将“一个零点可作为另一个零点的地址”理解为\emph{有限前缀的确定性寻址}的命题都不良定。
\par
在本文多轴前缀打包 \(a=(a_{\mathrm{vis}},a_{\mathrm{res}},a_{\mathrm{mode}})\)（定义 \ref{def:multi_axis_address}）中，上述结论可被解释为：当可见轴 \(a_{\mathrm{vis}}\) 仅取相位前缀且深度不足时，残差轴 \(a_{\mathrm{res}}\) 的最小预算被强制为 \(\Theta(2\log T)\) 级（否则必发生碰撞型 \(\mathsf{NULL}\)）；这与离线方向的端点二阶径向通道门槛同阶（推论 \ref{cor:xi-null-second-order-radial-channel}）。
\end{remark}

\begin{remark}[解释性叙述（非推理前提）：强混合极限下跨窗寻址的结构否定]\label{rem:xi-phase-mixing-no-addressing}
若将归一化零点过程（以单位平均间距归一化）假设为在有限维分布意义下趋于某个平移不变且混合的极限点过程（例如 \(\operatorname{Sine}_2\) 或更一般的 \(\operatorname{Sine}_\beta\) 过程；参见 \cite{AriasDeReynaRodgers2024ConvergencePointsZetaZeros,ChhaibiNajnudel2018RigiditySineBeta} 的相关极限框架），则任意仅依赖有限局部窗口配置的“地址函数”对远距窗口配置的条件信息增益在距离 \(R\to\infty\) 时消失。
该混合性语义上对应：跨尺度、跨窗口的确定性“零点互寻址”与强随机混合范式不相容；能稳定存在的仅是由有限分辨率前缀诱导的局部聚类（命题 \ref{prop:xi-phase-adjacent-prefix-sharing}）以及由端点通道门槛诱导的可见/不可见分层（定理 \ref{thm:app-jensen-offcritical-kink-localization} 与推论 \ref{cor:xi-null-second-order-radial-channel}）。
\end{remark}

\endinput

\input{sections/body/zeta_finite_part/xi/subsubsec__xi-addressability-resolution-law}

\begin{remark}[Integer layer calibration lattice of unitary slice: root unit point and Frobenius defect]\label{rem:xi-null-frobenius-calibration}
Note \ref{rem:xi-null-finite-precision} ​​indicates that endpoint compactification will push the "infinite height" difference into the last bit of precision under limited precision, thereby inducing null/unverifiable values ​​at the implementation level.
To decouple such implementation errors from real spectral crossing evidence, a set of \textbf{discrete anchor points} can be introduced on the unitary slice as integer layer guardrails.
\par
Take the prime number \(p\) ​​and take the root unit point grid on the unit circle
\(
u=e^{2\pi i j/p}
\) (\(j=0,\dots,p-1\), satisfies \(u^p=1\)).
It is deduced from \ref{cor:witt-root-of-unity-tower-lock} that there is tower locking modulo \(p^k\) on the \(p\)-power tower:
\[
a_{p^k}(u)\equiv a_{p^{k-1}}(1)\pmod{p^k}.
\]
Furthermore, by the definition \ref{def:witt-frobenius-defect} ​​and the proposition \ref{prop:witt-frobenius-defect-divisibility}, Frobenius defective polynomial
\(
\Delta_{p,k}(u):=a_{p^k}(u)-a_{p^{k-1}}(u^p)
\)
satisfies the divisibility depth
\(
\Delta_{p,k}(u)\in p^k\ZZ[u]
\)。
\par
The above discrete anchor points give a "no \(\mathsf{NULL}\) ​​calibration grid": deviations incompatible with the divisibility/congruence guardrails at these root unit points should be prioritized as implementation or data pipeline failures (rather than spectral crossing, geometric \(\mathsf{NULL}\) or interpretive interpretation of statistical residuals); see also note \ref{rem:witt-root-of-unity-acceleration} for its use as an acceleration law.
\end{remark}

\begin{corollary}[The "second-order nature" of radial channels and the structural origin of \texorpdfstring{\(\mathsf{NULL}\)}{NULL}]\label{cor:xi-null-second-order-radial-channel}
Appendix \ref{app:unit-circle-phase-gate}’s diskization interface transfers the offline zero-point event of the completed \(\Xi_\zeta\) to an on-disk zero-point event: If there is a zero point \(\rho=\beta+\mathrm{i}\gamma\) that satisfies \(0<\beta<\tfrac12\), then the corresponding \(w_\rho\in\mathbb{D}\) makes \(F_\zeta(w_\rho)=0\) (theorem \ref{thm:app-rh-iff-disk-zero-free} and the proposition \ref{prop:app-offcritical-radius-compression}).
And under the area phase endpoint depth \(\delta_\gamma:=\tfrac12-\pi^{-1}\arctan|\gamma|\downarrow 0\), the radial depth satisfies the second-order compression law (corollary \ref{cor:app-offcritical-radius-second-order-radial-channel}):
\[
1-|w_\rho|^{2}=4\pi^{2}\Bigl(\tfrac12-\beta\Bigr)\delta_\gamma^{2}\bigl(1+o(1)\bigr).
\]
Therefore, in \emph{pure phase system} (unitary slice), the deviation of \(\beta-\tfrac12\) ​​does not appear as a first-order visible phase disturbance, but is structurally pressed into the second-order radial channel;
If the closure refuses to introduce any radial record axis (\(|u|\neq 1\) ​​or equivalently \(|w|\neq 1\)), then the channel is erased by the type system and the offline zero point event can only be degraded to "unaddressable" (\(\mathsf{NULL}\)) in the visible readout.
On the contrary, if the external radial register is pressed "Extended Save" (see \(\operatorname{proj}^{\mathrm e}\) ​​of note \ref{rem:xi-soundness-vs-completeness}), the above-mentioned second-order channel can be restored to the non-degenerate extended readout of "phase readout + radial residual".
\end{corollary}


\begin{remark}[Auditable alternative: Address completeness adequacy criterion (ACC-Rank) given by Hankel-rank stability]\label{rem:xi-acc-rank}
The so-called "completeness" gap can be expressed as: there are zero-point events that are not addressable under the \(\mathbf{PW}\) ​​closure of this paper (the readout is not defined, that is, \(\mathsf{NULL}\)), so soundness alone is not enough to conclude that "all zero-points are captured in the visible domain".
In the bridge package implementation scope of this paper, this gap can be compiled into a back-end stability criterion of \emph{regression testable}: along the \(\theta\)-scan of the unitary slice, if the limit nuclear family does not need to introduce new visible degrees of freedom in the sense of minimum realization, then any zero-point event must have a \(\mathbf{PW}\) address, so that it will not be read as \(\mathsf{NULL}\).
\par
More specifically, for each finite level \(\ell\) ​​and each phase sampling point \(\theta\), obtain the trace sequence of the normalized nuclear family \(U_{L,\ell}(\theta)\) used in the proposition \ref{prop:xi-unitary-slice-spectral-crossing}
\(
a^{(\ell)}_n(\theta):=\Tr(U_{L,\ell}(\theta)^n)
\) ​​(or equivalent resolvent-trace fragment), and in Hankel rank/minimum implementation dimension
\[
d^{(\ell)}_{\min}(\theta):=\mathrm{rank}\,H\!\bigl(a^{(\ell)}_\bullet(\theta)\bigr)
\]
(Theorem \ref{thm:hankel-rank-minimal}) measures its “visible degrees of freedom”.
If there exists a constant \(d_\ast\) ​​such that \(d^{(\ell)}_{\min}(\theta)\le d_\ast\) exists for all sufficiently large levels \(\ell\) and all phase samples \(\theta\), and this bound remains stable when encrypting the \(\theta\) grid and raising levels, then the address completeness of "zero point \(\neq\mathsf{NULL}\)" can be considered to be satisfied by certification.
\par
This route replaces completeness from a pure existence statement to an auditable "dimensional stability" check; for Hankel-rank's implementation scope and script organization, see note \ref{rem:two_layer_bridge_package_completed_minH}, and refer to table \ref{tab:fold_collision_moment_hankel_rank} (script \texttt{scripts/exp\_fold\_collision\_moment\_hankel\_rank.py}) as the baseline template for the backend dimension certificate.
\par
From the perspective of \(\mathsf{NULL}\) semantics, the above check corresponds to a gap that is orthogonal to "geometric/branching/statistical": if the visible audit is limited to the exchange graph of trace/primitive in a fixed finite-dimensional state space, then Hankel rank (theorem \ref{thm:hankel-rank-minimal}) gives a hard lower bound for \emph{minimum record axis dimension}; any attempt to reduce the implementation dimension to \(<\mathrm{rank}(H)\) All strategies will cause irremovable ambiguity on certain visible observations, and the semantic consequence is "dimensional type \(\mathsf{NULL}\)" (the address must be externalized to a higher dimensional axis).
The gap can be certified in polynomial time by the Hankel normal form \(A=(H_0^{(\ell)})^{-1}H_1^{(\ell)}\) ​​(note \ref{rem:pom-s2-hankel}) and aligned with the stability test of the \(\theta\)-scan.
When further assuming "unified minimality + protocol closure (no unrecorded axis backflow)", the above ACC-Rank stability check can be upgraded from a sufficient criterion to a sufficient and necessary characterization, see theorem \ref{thm:xi-completeness-iff-hankel-bounded}.
\end{remark}

\begin{theorem}[Model axis completeness \texorpdfstring{$\Longleftrightarrow$}{<->} Hankel rank-uniform bounded (under unified minimum and protocol closure)]\label{thm:xi-completeness-iff-hankel-bounded}
Follow the settings and notation noted in \ref{rem:xi-acc-rank}, and consider the mode axis \(a_{\mathrm{mode}}\) ​​in the multi-axis prefix (definition \ref{def:multi_axis_address}).
Fix a confidence level so that the budget of the residual axis \(a_{\mathrm{res}}\) ​​is enough to exclude collision nulls caused by "identical collisions/label out-of-bounds" (such as satisfying the prime-register quantile threshold condition of proposition \ref{prop:xi-null-three-way}, or equivalently satisfying clause (2) of proposition \ref{prop:xi-completeness-auditable-decomposition}).
\par
Suppose further:
\begin{enumerate}
  \item \textbf{Uniform Minimality:} For all sufficiently large levels \(\ell\) ​​and all phase parameters \(\theta\), the “required visible degrees of freedom” of the trace sequence \(a^{(\ell)}_\bullet(\theta)\) is uniquely characterized by its Hankel rank, that is
  \[
  d^{(\ell)}_{\min}(\theta)=\mathrm{rank}\,H\!\bigl(a^{(\ell)}_\bullet(\theta)\bigr)
  \]
  And this minimum dimension is taken to be a lower bound on the mode axis capacity within the protocol closure that must be explicitly externalized (Theorem \ref{thm:hankel-rank-minimal}).
  \item \textbf{Protocol closure (no undocumented axes reflow):} Any added degrees of freedom affecting visible readouts/certificates that are not explicitly recorded as mode axes/registers are rejected at the type level and read out as \(\mathsf{NULL}\) ​​(defining \ref{def:xi-visible-read-null}; see also the "no undocumented axes" principle of theorem \ref{thm:xi-pw-no-continuous-hair}).
\end{enumerate}
Then the following statements are equivalent:
\begin{enumerate}
  \item[(i)] \textbf{Mode axis completeness:} There exists a constant \(d_\ast\) and a level threshold \(\ell_0\) such that all \(\ell\ge \ell_0\) and all phase samples \(\theta\), unitary-slice The visible degrees of freedom required for auditing can all be closed within the same finite-dimensional mode type; equivalently, zero-point events do not degenerate into \(\mathsf{NULL}\) under the \(\mathbf{PW}\) closure simply because "new visible degrees of freedom/mode axes need to be introduced".
  \item[(ii)] \textbf{Hankel rank-uniform bounded:}There are constants \(d_\ast\) ​​and \(\ell_0\) such that for all \(\ell\ge \ell_0\) and all \(\theta\)
  \[
  d^{(\ell)}_{\min}(\theta)\le d_\ast.
  \]
\end{enumerate}
\end{theorem}

\begin{proof}
\textbf{(ii)\(\Rightarrow\)(i).}
If there is \(d_\ast,\ell_0\) ​​such that \(d^{(\ell)}_{\min}(\theta)\le d_\ast\) is true for all \(\ell\ge\ell_0\) and \(\theta\), then for each \((\ell,\theta)\) there is a linear realization whose dimension does not exceed \(d_\ast\), realizing all the visible information of the trace sequence \(a^{(\ell)}_\bullet(\theta)\) (theorem \ref{thm:hankel-rank-minimal}).
By treating this realization (or its zero-padding extension on dimension \(d_\ast\)) as a unified mode axis type, the corresponding certificate readout can be completed without introducing new visible degrees of freedom; as assumed by the protocol closure, any object requiring additional degrees of freedom must be explicitly externalized, otherwise \(\mathsf{NULL}\) will be triggered, so that no new dimensional type \(\mathsf{NULL}\) will appear under the current "unified finite dimension" system.
Therefore, mode axis completeness holds.
\par
\textbf{(i)\(\Rightarrow\)(ii).}
Suppose Hankel rank inconsistency is bounded, then for any candidate constant \(d\) ​​and threshold \(\ell_0\) we can find \(\ell\ge\ell_0\) and a certain \(\theta\) such that
\(
d^{(\ell)}_{\min}(\theta)>\!d
\)。
By the unified minimality and theorem \ref{thm:hankel-rank-minimal}, the dimension of any linear representation that realizes this trace sequence is at least \(d^{(\ell)}_{\min}(\theta)\), so that irreducible modal collapse must occur within any fixed \(d\)-dimensional pattern type.
By protocol closure, "undocumented axis reflow" is prohibited: the unverifiable ambiguity induced by this collapse cannot be implicitly absorbed within the closure, and can only be triggered by mode axis deficiencies \(\mathsf{NULL}\).
This contradicts the requirement of (i) that "there exists a unified finite-dimensional pattern type such that events do not degenerate into \(\mathsf{NULL}\) ​​due to missing axes."
Therefore, the Hankel rank must be uniformly bounded, and we get (ii).
\end{proof}

\begin{proposition}[Auditable decomposition of completeness: Dimensionality Stability Certificate \(\times\) ​​Quota Certificate]\label{prop:xi-completeness-auditable-decomposition}
Under the unitary-slice zero-point audit scope of this paper, the structural sources of "zero-point event read as \(\mathsf{NULL}\)" can be classified into two categories: \emph{extra visible degrees of freedom} (protocol null value) and \emph{same image collision/label out-of-bounds} (collision null value, proposition \ref{prop:xi-null-three-way}).
Therefore, if the following two auditable constraints are met at the same time:
\begin{enumerate}
  \item \textbf{Visible degree of freedom stability (ACC-Rank):} There exists a constant \(d_\ast\) such that \(d^{(\ell)}_{\min}(\theta)\le d_\ast\) exists for all sufficiently large levels \(\ell\) and all phase samples \(\theta\), and this bound remains stable when densifying the \(\theta\) grid and increasing the level (Note \ref{rem:xi-acc-rank}; and see the theorem \ref{thm:xi-completeness-iff-hankel-bounded}) under unified minimality and protocol closure);
  \item \textbf{Adequate side information budget (quantile quota):} For any given \(\varepsilon\in(0,1)\), there is a set of prime-register modular product budget \(P_m\) such that
  \(
  P_m>2B_m(\varepsilon)=2\exp(m\gamma_m(\varepsilon))
  \)
  (Definition \ref{def:pom-gamma-quantile}), so the failure probability of "label out-of-bounds leading to non-unique reconstruction" can be reduced to \(\le\varepsilon\) ​​(Corollary \ref{cor:pom-prime-register-quantile-budget});
\end{enumerate}
Then under this confidence scope, the two unique sources (degree of freedom instability or tag collision) of the zero-point event reading \(\mathsf{NULL}\) are eliminated, so that "address completeness" can be guaranteed alternatively by the above two types of certificates in the system; merged with the theorem \ref{thm:xi-offset-null-type-safety} (soundness), that is, the completeness required by \(\Xi\)-RH Gap compression is \emph{Verification tasks for two types of backend certificates}: dimensionally stable certificates (ACC-Rank) and prime-register quota certificates (quantile threshold).
\end{proposition}

\begin{remark}[Counterexample collapse template: dual checker rejection in radial mode]\label{rem:xi-radial-collapse-template}
Section \ref{subsec:xi-radial-counterexample-collapse} gives a pseudo-object that is still in the form of a self-dual completion but enforces \(|u|\neq 1\) radial mode, and shows how it works in the root interval checker (proposition \ref{prop:xi-rh-interval-criterion}) and the constant level strictification checker (definition \ref{def:pom-finite-part-counterterm-gate}; proposition \ref{prop:pom-counterterm-anom-cancel}) failed at the same time, thereby explicitly turning the rejection chain of "untypeable \(\Rightarrow\ \mathsf{NULL}\)" into a reviewable audit process.
\end{remark}

\begin{remark}[Unified sealing: \texorpdfstring{$\mathsf{NULL}$}{NULL} under unitary slice scope]\label{rem:xi-null-conservation-closure}
In unitary slice semantics, "information conservation" is not equivalent to "\(\mathsf{NULL}\) ​​never occurs";
It is equivalent to compiling every source of nulls into an auditable ledger of gate conditions and defects, with external orthogonal axes of record where necessary to restore reversibility.
Under the scope of this paper, it can be summarized as the following closed loop:
\begin{enumerate}
  \item \textbf{Rejected outside the geometric domain.} If \(|u|\neq 1\) ​​(that is, \(u\notin\mathsf{Phase}\)), it is directly read as \(\mathsf{NULL}\) (definition \ref{def:xi-visible-read-null}) according to the visible domain constraints.
  \item \textbf{Defect quantification of intra-domain ambiguity.} When the macro edge \(\pi\) ​​is fixed, the invisible window width of the lift is accurately quantified by \(\kappa_m(\pi)\); the fiber shape residual of any lift is carried by \(D_{\mathrm{KL}}(\mu\|\pi^{\mathrm{e}})\) (theorem \ref{thm:xi-null-fiber-entropy-window} and corollary \ref{cor:xi-null-statistical-radius}).
  \item \textbf{Dimension threshold of finite-dimensional closure.} If the visible audit requires closure in a fixed finite-dimensional state space, then Hankel rank gives a lower bound on the minimum realized dimension; violation of this lower bound triggers the dimensionality type \(\mathsf{NULL}\) ​​(note \ref{rem:xi-acc-rank}).
  \item \textbf{Priority guardrails for the integer-level calibration lattice.} The Frobenius defect-divisible depth at the root unit point provides a discrete anchor point so that a large number of spurious \(\mathsf{NULL}\) ​​caused by finite precision can be prioritized as implementation failures at the integer level (note \ref{rem:xi-null-frobenius-calibration}).
\end{enumerate}
\end{remark}

\begin{definition}[Complete address as multi-axis prefix]\label{def:multi_axis_address}
Under the visible reading convention of this paper (definition \ref{def:xi-visible-read-null}) and the "address before value" semantics, a checkable and referable address should be understood as a typed \textbf{multi-axis prefix}
\[
a=\bigl(a_{\mathrm{vis}},a_{\mathrm{res}},a_{\mathrm{mode}}\bigr).
\]
The three axes respectively play the following roles:
\begin{enumerate}
  \item \textbf{Visible axis (\(a_{\mathrm{vis}}\))}: Select a visible column prefix/certificate prefix (e.g. the phase parameter \(\theta\) on the unitary slice and its induced auditable events) in the filtering of the given visible domain \(X_{\mathrm{vis}}\) with the checker \(\mathrm{Chk}\), thereby implementing the "concept" into a verifiable projection readout.
  \item \textbf{Residual axis (\(a_{\mathrm{res}}\))}: The record axis/tag axis used to injective the possible many-to-one compressed readout (Proposition \ref{prop:xi-null-three-way}). In the folded readout \(\Fold_m\) scenario, the worst-case capacity and average information budget of this axis are given hard lower bounds by the fiber multiplicity \(d_m\) and the defect scalar \(\kappa_m(\pi)\) respectively (theorem \ref{thm:address-defect-law}).
  \item \textbf{Mode axis (\(a_{\mathrm{mode}}\))}: Other orthogonal degrees of freedom required by the protocol contract (if not external, they are prohibited at the type level and read as \(\mathsf{NULL}\)). Typical examples include: radial register/amplitude channel (corollary \ref{cor:xi-null-second-order-radial-channel}) and minimum implementation dimension/visible degree of freedom stable axis (note \ref{rem:xi-acc-rank}).
\end{enumerate}
\par
Under this package, "whether the address is legal (not \(\mathsf{NULL}\))" means that the three axes meet at the same time: \(a_{\mathrm{vis}}\) falls within the visible domain filter, \(a_{\mathrm{res}}\)'s budget is enough to avoid co-image collision, and \(a_{\mathrm{mode}}\) is not prohibited by the protocol; any axis is missing or prohibited will trigger \(\mathsf{NULL}\) from the corresponding source, thus materializing the "completeness gap" into an auditable axis budget/axis stability condition (proposition \ref{prop:xi-completeness-auditable-decomposition}).
\end{definition}

\begin{remark}[Dual swap symmetry and residual budget obstacles]\label{rem:xi-dual-swap-symmetry-defect}
The "three-axis prefix" that defines \ref{def:multi_axis_address} describes the direct sum of the interface: \(a_{\mathrm{vis}}\) ​​is used as the quotient coordinate in the visible domain, and \(a_{\mathrm{res}}\) is used as the external label of fiber collision. The "orthogonality" of the two is structural decomposition rather than geometric similarity.
Therefore, the statement "two orthogonal axes are equivalent" in the semantics of this paper can only be understood as a stronger \textbf{duality interchangeable symmetry} event: there is an automorphism allowed by the protocol to interchange "visible quotient coordinates" with "fiber labels" and keep the checkable semantics unchanged.
\par
This swap generally does not hold, and its failure is given a strict obstacle by the residual budget.
Take the folded readout \(\Fold_m\) as an example: if there is \(x\in X_m\) that makes the fiber multiplicity \(d_m(x)=\card{\Fold_m^{-1}(x)}>1\), then any residual axis that makes the extended readout \((\Fold_m,r)\) an injective must provide at least \(d_m(x)\) distinguishable labels on the fiber (theorem \ref{thm:address-defect-law} Item (1)); at a macro scale, a lower bound on the average information budget must be paid \(H(r(A)\mid X)\ge \kappa_m(\pi)\) (Same as theorem (2)).
Thus, unless the fiber defects disappear on the subclass of interest (e.g. \(d_m(x)=1\) ​​or the corresponding \(\kappa_m(\pi)=0\)), the swap of "pressing fiber coordinates back to visible coordinates" cannot occur without externalizing the residual axis.
\end{remark}
